\documentclass{beamer}
\usetheme{Madrid}
\usecolortheme{default}
\usepackage[utf-8]{inputenc}
\usepackage[french]{babel}
\usepackage{tikz}
\usepackage{hyperref}
\usepackage{graphicx}
\usepackage{booktabs}

% Custom colors
\definecolor{RurDynamixGreen}{RGB}{76, 175, 80}
\definecolor{RurDynamixBlue}{RGB}{33, 150, 243}
\definecolor{DarkGray}{RGB}{64, 64, 64}

\setbeamercolor{primary}{fg=white,bg=RurDynamixGreen}
\setbeamercolor{secondary}{fg=white,bg=RurDynamixBlue}
\setbeamertemplate{frametitle}{
  \begin{beamercolorbox}[rounded=true,shadow=true,leftskip=1cm,rightskip=1cm,wd=\paperwidth,ht=0.8cm,dp=0.4cm]{frametitle}
    \usebeamerfont{frametitle}\insertframetitle
  \end{beamercolorbox}
}

% Title Page Setup
\title{Bassins artificiels et dynamiques agro-environnementales \\ dans le Rif marocain}
\subtitle{Analyse spatio-temporelle par télédétection et SIG (2010-2025)}
\author{Ousmane Samba MAÏGA \\ \small{Sous la direction de : Prof. Said Moukrim}}
\institute{Rur\_Dynamix Conference 2026 \\ Rural Development and Environmental Dynamics}
\date{\today}

\begin{document}

% Slide 1: Title
\frame{\titlepage}

% Slide 2: Contexte et Problématique
\begin{frame}
\frametitle{Contexte et Problématique}
\begin{itemize}
  \item \textbf{Transformation rapide du Rif central} : Expansion massive de la culture du cannabis depuis la légalisation (2021)
  \item \textbf{Innovations d'adaptation} : Prolifération de bassins artificiels pour l'irrigation (25 → 5,139 bassins, 2010-2025)
  \item \textbf{Questions critiques} :
  \begin{itemize}
    \item Quels facteurs biophysiques expliquent l'implantation des bassins ?
    \item Quel est leur impact réel sur le couvert végétal (NDVI) ?
    \item Quelles implications pour la durabilité ?
  \end{itemize}
  \item \textbf{Urgence} : Changement climatique + intensification agricole → Risques hydriques et environnementaux
\end{itemize}
\end{frame}

% Slide 3: Objectifs de l'étude
\begin{frame}
\frametitle{Objectifs de l'étude}
\textbf{Objectif général :} Analyser la dynamique spatio-temporelle des bassins artificiels dans le Rif et évaluer leurs impacts sur le couvert végétal
\\[0.5cm]
\textbf{Objectifs spécifiques :}
\begin{enumerate}
  \item Cartographier et quantifier l'évolution des bassins (2010-2025)
  \item Analyser les facteurs de localisation (altitude, pente, hydrographie)
  \item Évaluer l'impact sur le couvert végétal (NDVI, 1998-2026)
  \item Identifier les patterns de densité des bassins et leurs effets
\end{enumerate}
\end{frame}

% Slide 4: Zone d'étude
\begin{frame}
\frametitle{Zone d'étude}
\begin{columns}[T]
  \column{0.5\textwidth}
  \textbf{Localisation :}
  \begin{itemize}
    \item Province de Chefchaouen (Rif central)
    \item Communes : Bab Berred, Bni Darkoul, Bni Salah, Oued Malha
    \item Altitude : 300-2,000 m
    \item Relief montagneux, pentes fortes
  \end{itemize}

  \column{0.5\textwidth}
  \textbf{Caractéristiques :}
  \begin{itemize}
    \item Climat : Semi-aride à subhumide
    \item Agriculture intensive (cannabis légal)
    \item Ressources hydriques limitées
    \item Écosystème fragile en Méditerranée
  \end{itemize}
\end{columns}

\begin{center}
  \small{\textit{Figure 1 : Carte de localisation}}
\end{center}
\end{frame}

% Slide 5: Méthodologie
\begin{frame}
\frametitle{Méthodologie}
\textbf{Données et outils :}
\begin{itemize}
  \item \textbf{Télédétection :} Images Landsat (1998-2026), Google Earth, Sentinel-2
  \item \textbf{SIG :} QGIS, ArcGIS pour cartographie et analyse spatiale
  \item \textbf{Traitement :} RStudio, Google Earth Engine (GEE)
  \item \textbf{Indices :} NDVI pour suivi du couvert végétal
\end{itemize}

\textbf{Approche analytique :}
\begin{enumerate}
  \item Digitalisation manuelle des bassins (2010-2025)
  \item Analyse morphométrique (surface, volume, profondeur)
  \item Analyse de localisation (altitude, pente, proximité réseau hydro)
  \item Suivi NDVI multi-temporel + corrélation précipitations
  \item Analyse de densité et effets spatiaux
\end{enumerate}
\end{frame}

% Slide 6: Résultats - Dynamique des bassins
\begin{frame}
\frametitle{Résultats - Explosion du nombre de bassins}
\begin{center}
  \textbf{Croissance exceptionnelle : 25 → 5,139 bassins (2010-2025)}
  \\[0.3cm]
  \small{Augmentation de plus de 20,000\%}
\end{center}

\begin{itemize}
  \item \textbf{Période d'accélération} : 2013-2025 (légalisation du cannabis, 2021)
  \item \textbf{Miniaturisation} : Bassins de plus en plus petits et nombreux
  \item \textbf{Distribution inégale} : Concentration dans communes spécifiques
  \item \textbf{Anarchie spatiale} : Absence de planification centralisée
\end{itemize}

\begin{center}
  \small{\textit{Figure 2 : Évolution du nombre cumulé de bassins (2010-2025)}}
\end{center}
\end{frame}

% Slide 7: Résultats - Facteurs de localisation
\begin{frame}
\frametitle{Résultats - Facteurs de localisation}
\textbf{L'implantation répond à des contraintes biophysiques :}
\begin{columns}[T]
  \column{0.5\textwidth}
  \textbf{Altitude} :
  \begin{itemize}
    \item Concentration 600-1,200 m
    \item Évite zones très élevées
  \end{itemize}

  \textbf{Pente} :
  \begin{itemize}
    \item Préférence : 15-35\% (eau coule)
    \item Adaptation au relief
  \end{itemize}

  \column{0.5\textwidth}
  \textbf{Réseau hydrographique} :
  \begin{itemize}
    \item Proximité critique (< 500m)
    \item Accès à l'eau déterminant
  \end{itemize}

  \textbf{Occupation du sol} :
  \begin{itemize}
    \item Préférence : zones agricoles
    \item Conversion de forêt
  \end{itemize}
\end{columns}

\begin{center}
  \small{\textit{Figures 3-4 : Cartes de répartition par altitude, pente, hydrographie}}
\end{center}
\end{frame}

% Slide 8: Résultats - Impact sur NDVI
\begin{frame}
\frametitle{Résultats - Impact sur le couvert végétal (NDVI)}
\textbf{Paradoxe détecté (1998-2026) :}
\begin{itemize}
  \item \textbf{Période 1998-2013} : NDVI stable/légèrement décroissant
  \item \textbf{Période 2013-2026} : Amélioration significative du couvert végétal
  \item \textbf{Rupture} : 2013 coïncide avec essor des bassins
  \item \textbf{Corrélation} : Forte avec irrigation (bassins ?) mais aussi avec précipitations
\end{itemize}

\textbf{Interprétation :}
\begin{itemize}
  \item Irrigation via bassins → amélioration du NDVI ?
  \item Mais aussi : évolution climatique, cycle de précipitations
  \item Pas de causalité directement établie
  \item Besoin d'étude plus fine (parcelle-échelle)
\end{itemize}

\begin{center}
  \small{\textit{Figures 5-6 : Cartes NDVI comparatif, corrélation précipitations}}
\end{center}
\end{frame}

% Slide 9: Résultats - Synthèse quantitative
\begin{frame}
\frametitle{Résultats - Synthèse : Morphométrie des bassins}
\begin{center}
  \small{
  \begin{tabular}{|c|c|c|c|}
    \hline
    \textbf{Paramètre} & \textbf{Min} & \textbf{Max} & \textbf{Moyen} \\
    \hline
    Superficie (m²) & 50 & 50,000 & 2,500 \\
    Profondeur (m) & 1 & 8 & 3 \\
    Volume (m³) & 100 & 300,000 & 7,500 \\
    \hline
  \end{tabular}
  }
\end{center}

\textbf{Tendances :}
\begin{itemize}
  \item Miniaturisation progressive des bassins
  \item Bassin moyen : 2,500 m² (récemment < 1,000 m²)
  \item Volume total estimé : ~38 millions m³ (toute zone)
\end{itemize}

\begin{center}
  \small{\textit{Figure 7 : Distribution de superficie/volume/profondeur}}
\end{center}
\end{frame}

% Slide 10: Discussion et enjeux
\begin{frame}
\frametitle{Discussion - Enjeux et défis}
\textbf{Questions de durabilité :}
\begin{itemize}
  \item \textbf{Hydrologie} : Prélèvement excessif ? Dégradation nappes phréatiques ?
  \item \textbf{Environnemental} : Perte de forêt (5,139 bassins = impact cumulatif)
  \item \textbf{Social} : Inégalités d'accès à l'eau, conflits (cannabis vs vivrières)
  \item \textbf{Adaptation} : Efficacité face au climat changeant ? Résilience ?
\end{itemize}

\textbf{Lacunes de connaissance :}
\begin{itemize}
  \item Lien précis NDVI-bassins (causalité vs corrélation)
  \item Impact hydrologique (débit réseau hydro, infiltration)
  \item Efficacité énergétique et économique
\end{itemize}
\end{frame}

% Slide 11: Recommandations
\begin{frame}
\frametitle{Recommandations et perspectives}
\textbf{Pour la recherche :}
\begin{itemize}
  \item Étude hydrologique fine : bilans hydrologiques, débit, prélèvement
  \item Modélisation : scénarios futurs (climat, cannabis, eau)
  \item Approches participatives : perception des agriculteurs
\end{itemize}

\textbf{Pour la gestion territoriale :}
\begin{itemize}
  \item Planification centralisée de l'irrigation
  \item Suivi continu par télédétection (alerte précoce)
  \item Régulation de la taille/densité des bassins
  \item Adaptation aux conditions hydriques changeantes
\end{itemize}

\textbf{Pour la politique :}
\begin{itemize}
  \item Intégrer eau et agriculture dans stratégie climatique
  \item Équilibre cannabis-cultures vivrières-environnement
\end{itemize}
\end{frame}

% Slide 12: Conclusion
\begin{frame}
\frametitle{Conclusion}
\begin{center}
  \textbf{\Large{Synthèse et perspectives}}
\end{center}

\textbf{Principaux messages :}
\begin{itemize}
  \item Le Rif central connaît une transformation rapide pilotée par l'expansion du cannabis et l'innovation hydraulique
  \item Bassin artificiels = adaptation ingénieuse aux contraintes hydriques locales
  \item Mais : impacts cumulatifs mal quantifiés → urgence de suivi/régulation
  \item Paradoxe NDVI → amélioration du couvert ?  Mais au prix de ressources ?
\end{itemize}

\textbf{Enjeu central :}
\begin{center}
  \textit{Concilier productivité agricole, résilience climatique \\ et durabilité environnementale en contexte Méditerranéen}
\end{center}
\end{frame}

% Final Slide
\begin{frame}[plain]
\begin{center}
  \vfill
  {\Huge \textbf{Merci de votre attention}}
  \\[1cm]
  {\Large Questions ?}
  \\[2cm]
  \small{
    Ousmane Samba MAÏGA \\
    Master BioEnv2D \\
    Faculté des Sciences de Rabat
  }
  \vfill
\end{center}
\end{frame}

\end{document}
